\section{Corrections and Clarifications}
\label{sec:source-corrections}

Surprisingly but excitingly, we found several errors in the original paper. They do not overturn the main conclusions, but should not be ignored.

The probability statement for Algorithm~2 requires additional accounting. Let
\[
  p_n=\left(1-n^{-1/8}\right)\left(1-\frac{2}{n}\right)^2.
\]
Lemma \ref{lem:query-efficient-eps-wsne-in-zero_sum-games} guarantees success probability at least $p_n$ for \textbf{one} invocation of the zero-sum routine, whereas Algorithm~2 invokes it \textbf{twice}. With independent random coins, both invocations therefore succeed with probability at least $p_n^2$.

Moreover, The paper does not account for \textbf{sampling failures}. Lemma \ref{lem:small-supp-approx} is existential, but a randomized implementation must draw the small-support profile and include the sampling event. Hoeffding bounds and a union bound over its constraints show that
\[
  s=O\!\left(\frac{\log n+\log(1/\eta)}{\varepsilon^2}\right)
\]
samples suffice for failure probability at most $\eta$ \cite{hoeffding63}. Consequently, a complete success guarantee is $p_n^2(1-\eta)$; for example, taking $\eta=1/n$ gives
\[
  \left(1-n^{-1/8}\right)^2
  \left(1-\frac{2}{n}\right)^4
  \left(1-\frac1n\right).
\]
This correction does not change the approximation guarantee. The query count becomes
\[
  O\!\left(
  \frac{n\log n}{\varepsilon^4}
  +\frac{n(\log n+\log(1/\eta))}{\varepsilon^2}
  \right),
\]
which remains $O(n\log n/\varepsilon^4)$ for $\eta=1/n$. Algorithm~1 enumerates at most $O(n^{2\kappa(\delta)})$ profiles, while Algorithm~2 enumerates at most $O(s^{\kappa(\delta)}n^{\kappa(\delta)})$ profiles in its queried subgame; hence both running times remain polynomial in $n$ for fixed $\varepsilon$ and $\delta$. Different allocations of the failure probabilities produce the same asymptotic bounds.

The source also contains a few insignificant typos:
\begin{itemize}
  \item The introduction claims that $\varepsilon$-WSNEs in zero-sum bimatrix games can be found within $O\left(\frac{n \log n}{\varepsilon^2}\right)$ queries. But by Theorem~\ref{thm:query-efficient-version}, the algorithm can perform up to $O\left(\frac{n \log n}{\varepsilon^4}\right)$ queries.
  \item The proof of Theorem~\ref{thm:query-efficient-version} \cite[proof of Theorem~2]{12WSNE} invokes Lemma~\ref{lem:small-supp-approx} \cite[Lemma~1]{12WSNE} in Case~3(b), although the required $3\varepsilon$-WSNE conclusion is supplied by Lemma~\ref{lem:3_eps-wsne-samling} \cite[Lemma~2]{12WSNE}.
  \item In the proof of Theorem~\ref{thm:query-efficient-version} \cite[proof of Theorem~2]{12WSNE}, the first paragraph states $|\operatorname{supp}(\mathbf{x}_s^*)| = O\left(\frac{\log n}{\varepsilon^2}\right)$ twice. The second occurrence should be $|\operatorname{supp}(\mathbf{y}_s^*)| = O\left(\frac{\log n}{\varepsilon^2}\right)$.
\end{itemize}

